\documentclass{notaaula}

        \titulonota{Indução e ondas eletromagnéticas}
        \creditos{Turma 3C · Física · Novembro/2026 · Prof. Josué Cavalcante}

        \begin{document}
        \cabecalho

        \begin{sobre}
        Eletromagnetismo: campos produzidos por correntes, indução, motores, geradores, transformadores e ondas eletromagnéticas. Habilidades em foco: EM13CNT303 e EM13CNT308.
        \end{sobre}

        \section{Corrente produz campo magnético}

\begin{copiar}{Experiência de Oersted}
Uma corrente elétrica desvia uma bússola próxima: cargas em movimento produzem campo magnético.
Em fio longo e retilíneo,
\[
  B=\frac{\mu_0 i}{2\pi r},
  \qquad \mu_0=4\pi\times10^{-7}\un{T\,m/A}.
\]
As linhas são circunferências ao redor do fio. A regra da mão direita fornece seu sentido.
Em solenoide longo, aproximadamente,
\[
  B=\mu_0 n i,
\]
onde $n$ é o número de espiras por metro.
\end{copiar}

\begin{exemplo}
A $\dec{0,10}\un{m}$ de um fio que conduz $5\un{A}$,
\[
  B=\frac{4\pi\times10^{-7}\cdot5}{2\pi\cdot\dec{0,10}}
  =\resultado{1\times10^{-5}\un{T}}.
\]
\end{exemplo}

\section{Fluxo e indução}

\begin{copiar}{Faraday e Lenz}
O fluxo magnético por uma espira plana em campo uniforme é
\[
  \Phi=B A\cos\theta.
\]
Uma variação do fluxo induz força eletromotriz:
\[
  \mathcal E=-N\frac{\Delta\Phi}{\Delta t}.
\]
O sinal negativo representa a lei de Lenz: a corrente induzida cria efeito magnético que se opõe à
variação de fluxo que a produziu. Variação pode resultar de mudança em $B$, na área, no ângulo ou
do movimento relativo entre ímã e circuito.
\atencao{Fluxo constante não induz fem, mesmo que o fluxo seja grande. É necessário haver variação.}
\end{copiar}

\begin{exemplo}[Uma espira]
O fluxo muda de $\dec{0,08}\un{Wb}$ para $\dec{0,02}\un{Wb}$ em $\dec{0,03}\un{s}$.
O módulo da fem média é
\[
  |\mathcal E|=\frac{|\Delta\Phi|}{\Delta t}
  =\frac{\dec{0,06}}{\dec{0,03}}
  =\resultado{2\un{V}}.
\]
\end{exemplo}

\section{Motores, geradores e transformadores}

\begin{copiar}{Conversões de energia}
\begin{itens}
\item Motor: energia elétrica $\rightarrow$ energia mecânica, usando força magnética sobre correntes.
\item Gerador: energia mecânica $\rightarrow$ energia elétrica, usando indução.
\item Transformador: altera tensão e corrente alternadas por indução entre bobinas.
\end{itens}
Para transformador ideal,
\[
  \frac{U_s}{U_p}=\frac{N_s}{N_p},
  \qquad U_pi_p=U_si_s.
\]
Transmitir energia em alta tensão reduz corrente e perdas $P_{\mathrm{perda}}=Ri^2$ nos cabos.
\end{copiar}

\section{Ondas eletromagnéticas}

\begin{copiar}{Espectro}
Campos elétrico e magnético variáveis podem se propagar como onda eletromagnética. No vácuo,
\[
  c=\lambda f.
\]
Da menor para a maior frequência: rádio, micro-ondas, infravermelho, visível, ultravioleta, raios X
e raios gama. Frequência maior significa fóton mais energético. Radiofrequência não é ionizante;
raios X e gama são ionizantes e exigem controle de dose.
\end{copiar}

\begin{dica}
Transformador não funciona com corrente contínua estacionária, porque ela não mantém fluxo variável
na bobina secundária.
\end{dica}

\begin{exercicios}
\nivel{1}{Conceitos}
\begin{questoes}
\item Que evidência a experiência de Oersted forneceu?
\item O que deve ocorrer com o fluxo para aparecer fem induzida?
\item Enuncie a lei de Lenz em palavras.
\item Diferencie motor e gerador pela conversão de energia.
\item Organize rádio, visível e raios X por frequência crescente.
\nivel{2}{Aplicação}
\item Se a distância a um fio longo dobra, o que acontece com $B$?
\item Uma bobina de 20 espiras tem variação de fluxo de $\dec{0,01}\un{Wb}$ em $\dec{0,10}\un{s}$. Ache $|\mathcal E|$.
\item Um transformador ideal tem $N_p=1000$, $N_s=100$ e $U_p=220\un{V}$. Ache $U_s$.
\item Uma onda de rádio tem $f=100\un{MHz}$. Use $c=3\times10^8\un{m/s}$ e ache $\lambda$.
\nivel{3}{Síntese}
\item Explique por que alta tensão reduz perdas na transmissão para a mesma potência entregue.
\item Aproximar o polo norte de um ímã de uma espira aumenta o fluxo. Qual é a função do campo induzido?
\item Compare radiação ionizante e não ionizante sem confundir frequência com intensidade.
\end{questoes}
\gabarito{2) o fluxo deve variar; 5) rádio, visível, raios X; 6) cai à metade;
7) $2\un{V}$; 8) $22\un{V}$; 9) $3\un{m}$.}
\end{exercicios}


\end{document}
